\documentclass{amsart}
% For dealing with references we use the comment environment
\usepackage{verbatim}
\newenvironment{reference}{\comment}{\endcomment}
%\newenvironment{reference}{}{}
\newenvironment{slogan}{\comment}{\endcomment}
\newenvironment{history}{\comment}{\endcomment}

% For commutative diagrams we use Xy-pic
\usepackage[all]{xy}

% We use 2cell for 2-commutative diagrams.
\xyoption{2cell}
\UseAllTwocells

% We use multicol for the list of chapters between chapters
\usepackage{multicol}

% This is generally recommended for better output
\usepackage{lmodern}
\usepackage[T1]{fontenc}

% For cross-file-references

% Package for hypertext links:
\usepackage{hyperref}

% For any local file, say "hello.tex" you want to link to please

% Theorem environments.
%
\theoremstyle{plain}
\newtheorem{theorem}[subsection]{Theorem}
\newtheorem{proposition}[subsection]{Proposition}
\newtheorem{lemma}[subsection]{Lemma}

\theoremstyle{definition}
\newtheorem{definition}[subsection]{Definition}
\newtheorem{example}[subsection]{Example}
\newtheorem{exercise}[subsection]{Exercise}
\newtheorem{situation}[subsection]{Situation}

\theoremstyle{remark}
\newtheorem{remark}[subsection]{Remark}
\newtheorem{remarks}[subsection]{Remarks}

\numberwithin{equation}{subsection}

% Macros
%
\def\lim{\mathop{\mathrm{lim}}\nolimits}
\def\colim{\mathop{\mathrm{colim}}\nolimits}
\def\Spec{\mathop{\mathrm{Spec}}}
\def\Hom{\mathop{\mathrm{Hom}}\nolimits}
\def\Ext{\mathop{\mathrm{Ext}}\nolimits}
\def\SheafHom{\mathop{\mathcal{H}\!\mathit{om}}\nolimits}
\def\SheafExt{\mathop{\mathcal{E}\!\mathit{xt}}\nolimits}
\def\Sch{\mathit{Sch}}
\def\Mor{\mathop{\mathrm{Mor}}\nolimits}
\def\Ob{\mathop{\mathrm{Ob}}\nolimits}
\def\Sh{\mathop{\mathit{Sh}}\nolimits}
\def\NL{\mathop{N\!L}\nolimits}
\def\CH{\mathop{\mathrm{CH}}\nolimits}
\def\proetale{{pro\text{-}\acute{e}tale}}
\def\etale{{\acute{e}tale}}
\def\QCoh{\mathit{QCoh}}
\def\Ker{\mathop{\mathrm{Ker}}}
\def\Im{\mathop{\mathrm{Im}}}
\def\Coker{\mathop{\mathrm{Coker}}}
\def\Coim{\mathop{\mathrm{Coim}}}

% Boxtimes
%
\DeclareMathSymbol{\boxtimes}{\mathbin}{AMSa}{"02}

%
% Macros for moduli stacks/spaces
%
\def\QCohstack{\mathcal{QC}\!\mathit{oh}}
\def\Cohstack{\mathcal{C}\!\mathit{oh}}
\def\Spacesstack{\mathcal{S}\!\mathit{paces}}
\def\Quotfunctor{\mathrm{Quot}}
\def\Hilbfunctor{\mathrm{Hilb}}
\def\Curvesstack{\mathcal{C}\!\mathit{urves}}
\def\Polarizedstack{\mathcal{P}\!\mathit{olarized}}
\def\Complexesstack{\mathcal{C}\!\mathit{omplexes}}
% \Pic is the operator that assigns to X its picard group, usage \Pic(X)
% \Picardstack_{X/B} denotes the Picard stack of X over B
% \Picardfunctor_{X/B} denotes the Picard functor of X over B
\def\Pic{\mathop{\mathrm{Pic}}\nolimits}
\def\Picardstack{\mathcal{P}\!\mathit{ic}}
\def\Picardfunctor{\mathrm{Pic}}
\def\Deformationcategory{\mathcal{D}\!\mathit{ef}}

\setlength{\emergencystretch}{3em}
\begin{document}
\title[Stacks Project, Tag 059M]{1308 --- Mittag-Leffler modules over arbitrary rings are exactly those whose tensor functor preserves product injections}
\maketitle
\noindent Copyright (C) 2005--2025 Johan de Jong. Stacks Project authors. Source: \href{https://stacks.math.columbia.edu/tag/059M}{Tag 059M}.

\noindent Editorial difficulty note (not author text). Difficulty H2: The reverse direction packages every possible finite-presentation tensor-kernel relation into one universal product element, kills it at a single finite stage, and extracts simultaneous domination for every test module. The forward direction requires uniform stabilization across an arbitrary product; this is a substantive finiteness characterization beyond qualification algebra..

\noindent Editorial provenance note (not author text). History: excerpt from revision \texttt{a0444\allowbreak 6e57e\allowbreak c1fbc\allowbreak 252a8\allowbreak 71afc\allowbreak ec775\allowbreak 2fb28\allowbreak 07b14}, algebra.tex, lines 21442--21531. Reference presentation was adapted; original-author context and core support blocks were retained or added. The mathematical wording is unchanged. Licensed under GNU FDL 1.2 or later; no invariant sections or cover texts. See ../licenses/COPYING. Context blocks reproduce author definitions and supporting results; they are not additional counted problems.

\begin{definition}
\label{definition-mittag-leffler-module}
Let $M$ be an $R$-module.  We say that $M$ is {\it Mittag-Leffler} if the
equivalent conditions of
Proposition \ref{proposition-ML-characterization}
hold.
\end{definition}

\begin{definition}
\label{definition-domination}
Let $f: M \to N$ and $g: M \to M'$ be maps of $R$-modules.
Then we say $g$ {\it dominates} $f$ if for any $R$-module $Q$, we have $\Ker(f
\otimes_R \text{id}_Q) \subset \Ker(g \otimes_R \text{id}_Q)$.
\end{definition}

\begin{definition}
\label{definition-universally-injective}
Let $f: M \to N$ be a map of $R$-modules.  Then $f$ is called
{\it universally injective} if for every $R$-module $Q$, the map $f
\otimes_R \text{id}_Q: M \otimes_R Q \to N \otimes_R Q$
is injective.  A sequence $0 \to M_1 \to M_2 \to M_3
\to 0$ of $R$-modules is called {\it universally exact} if it is exact
and $M_1 \to M_2$ is universally injective.
\end{definition}

\begin{proposition}
\label{proposition-ML-characterization}
Let $M$ be an $R$-module.  Let $(M_i, f_{ij})$ be a directed system of finitely
presented $R$-modules, indexed by $I$, such that $M = \colim M_i$.  Let
$f_i:
M_i \to M$ be the canonical map.  The following are equivalent:
\begin{enumerate}
\item For every finitely presented $R$-module $P$ and module map $f: P
\to M$, there exists a finitely presented $R$-module $Q$ and a module
map $g: P \to Q$ such that $g$ and $f$ dominate each other, i.e.,
$\Ker(f \otimes_R \text{id}_N) = \Ker(g \otimes_R \text{id}_N)$
for every $R$-module $N$.
\item For each $i \in I$, there exists $j \geq i$ such that $f_{ij}: M_i
\to M_j$ dominates $f_i: M_i \to M$.
\item For each $i \in I$, there exists $j \geq i$ such that $f_{ij}: M_i
\to M_j$ factors through $f_{ik}: M_i \to M_k$ for all $k \geq
i$.
\item For every $R$-module $N$, the inverse system
$(\Hom_R(M_i, N), \Hom_R(f_{ij}, N))$ is Mittag-Leffler.
\item For $N = \prod_{s \in I} M_s$, the inverse system
$(\Hom_R(M_i, N), \Hom_R(f_{ij}, N))$ is Mittag-Leffler.
\end{enumerate}
\end{proposition}

\begin{proof}
First we prove the equivalence of (1) and (2).  Suppose (1) holds and let $i
\in I$. Corresponding to the map $f_i: M_i \to M$, we can choose $g:
M_i \to Q$ as in (1).  Since $M_i$ and $Q$ are of finite presentation,
so is $\Coker(g)$.  Then by Lemma \ref{lemma-domination}, $f_i : M_i
\to M$ factors through $g: M_i \to Q$, say $f_i = h \circ g$
for some $h: Q \to M$.  Then since $Q$ is finitely presented, $h$
factors through $M_j \to M$ for some $j \geq i$, say $h = f_j \circ h'$
for some $h': Q \to M_j$.  In total we have a commutative diagram
$$
\xymatrix{
  &  M  & \\
M_i \ar[dr]_g \ar[ur]^{f_i} \ar[rr]^{f_{ij}} &
& M_j \ar[ul]_{f_j} \\
  & Q  \ar[ur]_{h'} &
}
$$
Thus $f_{ij}$ dominates $g$.  But $g$ dominates $f_i$, so $f_{ij}$ dominates
$f_i$.

\medskip\noindent
Conversely, suppose (2) holds.  Let $P$ be of finite presentation and $f: P
\to M$ a module map.  Then $f$ factors through $f_i: M_i \to M$
for some $i \in I$, say $f = f_i \circ g'$ for some $g': P \to M_i$.
Choose by (2) a $j \geq i$ such that $f_{ij}$ dominates $f_i$.  We have a
commutative diagram
$$
\xymatrix{
P \ar[d]_{g'} \ar[r]^{f}           & M  \\
M_i \ar[ur]^{f_i} \ar[r]_{f_{ij}} & M_j \ar[u]_{f_j}
}
$$
From the diagram and the fact that $f_{ij}$ dominates $f_i$, we find that $f$
and $f_{ij} \circ g'$ dominate each other.  Hence taking $g = f_{ij} \circ g' :
P \to M_j$ works.

\medskip\noindent
Next we prove (2) is equivalent to (3).  Let $i \in I$.  It is always true that
$f_i$ dominates $f_{ik}$ for $k \geq i$, since $f_i$ factors through
$f_{ik}$.  If (2) holds, choose $j \geq i$ such that $f_{ij}$ dominates
$f_i$.  Then since domination is a transitive relation, $f_{ij}$ dominates
$f_{ik}$ for $k \geq i$. All $M_i$ are of finite presentation, so
$\Coker(f_{ik})$ is of finite presentation for $k \geq i$.  By Lemma
\ref{lemma-domination}, $f_{ij}$ factors through $f_{ik}$ for all $k \geq i$.
Thus (2) implies (3).  On the other hand, if (3) holds then for any $R$-module
$N$, $f_{ij} \otimes_R \text{id}_N$ factors through $f_{ik}
\otimes_R \text{id}_N$ for $k \geq i$.  So $\Ker(f_{ik} \otimes_R
\text{id}_N) \subset \Ker(f_{ij} \otimes_R \text{id}_N)$ for $k
\geq i$.  But $\Ker(f_i \otimes_R \text{id}_N: M_i \otimes_R N
\to M \otimes_R N)$ is the union of $\Ker(f_{ik} \otimes_R
\text{id}_N)$ for $k \geq i$.  Thus $\Ker(f_i \otimes_R
\text{id}_N) \subset \Ker(f_{ij} \otimes_R \text{id}_N)$ for
any $R$-module $N$, which by definition means $f_{ij}$ dominates $f_i$.

\medskip\noindent
It is trivial that (3) implies (4) implies (5).  We show (5) implies (3).  Let
$N = \prod_{s \in I} M_s$. If (5) holds, then given $i \in I$ choose $j \geq i$
such that
$$
\Im( \Hom(M_j, N) \to  \Hom(M_i, N)) =
\Im( \Hom(M_k, N) \to  \Hom(M_i, N))
$$
for all $k \geq j$.  Passing the product over $s \in I$ outside of the
$\Hom$'s
and looking at the maps on each component of the product, this says
$$
\Im( \Hom(M_j, M_s) \to  \Hom(M_i, M_s)) =
\Im( \Hom(M_k, M_s) \to   \Hom(M_i, M_s))
$$
for all $k \geq j$ and $s \in I$.  Taking $s = j$ we have
$$
\Im( \Hom(M_j, M_j) \to  \Hom(M_i, M_j)) =
\Im( \Hom(M_k, M_j) \to   \Hom(M_i, M_j))
$$
for all $k \geq j$.  Since $f_{ij}$ is the image of
$\text{id} \in  \Hom(M_j, M_j)$ under
$\Hom(M_j, M_j) \to  \Hom(M_i, M_j)$,
this shows  that for any $k \geq j$ there is $h \in \Hom(M_k, M_j)$
such that $f_{ij} = h \circ f_{ik}$.  If $j \geq k$ then we can take
$h = f_{kj}$. Hence (3) holds.
\end{proof}


\begin{lemma}
\label{lemma-domination}
Let $f: M \to N$ and $g: M \to M'$ be maps of $R$-modules.
Suppose $\Coker(f)$ is of finite presentation.  Then $g$ dominates $f$ if
and
only if $g$ factors through $f$, i.e.\ there exists a module map $h: N
\to M'$ such that $g = h \circ f$.
\end{lemma}

\begin{proof}
Consider the pushout of $f$ and $g$ as in the statement of
Lemma \ref{lemma-domination-universally-injective}.
From the construction of the pushout it follows that
$\Coker(f') = \Coker(f)$, so $\Coker(f')$ is of finite
presentation.  Then by
Lemma \href{https://stacks.math.columbia.edu/tag/058L}{lemma universally exact split (Tag 058L)}, $f'$ is universally injective if and
only if
$$
0 \to M' \xrightarrow{f'} N' \to \Coker(f') \to 0
$$
splits. This is the case if and only if there is a map $h' : N' \to M'$
such that $h' \circ f' = \text{id}_{M'}$.  From the universal
property of the pushout, the existence of such an $h'$ is equivalent to $g$
factoring through $f$.
\end{proof}


\begin{lemma}
\label{lemma-domination-universally-injective}
Let $f : M \to N$ and $g : M \to M'$ be maps of $R$-modules.
Consider the pushout of $f$ and $g$,
$$
\xymatrix{
M  \ar[r]_f \ar[d]_g & N \ar[d]^{g'} \\
M' \ar[r]^{f'} & N'
}
$$
Then $g$ dominates $f$ if and only if $f'$ is universally injective.
\end{lemma}

\begin{proof}
Recall that $N'$ is $M' \oplus N$ modulo the submodule consisting of elements
$(g(x), -f(x))$ for $x \in M$.
From the construction of $N'$ we have a short exact sequence
$$
0 \to \Ker(f) \cap \Ker(g) \to \Ker(f) \to \Ker(f')
\to 0.
$$
Since tensoring commutes with taking pushouts, we have such a short exact
sequence
$$
0 \to \Ker(f \otimes \text{id}_Q ) \cap \Ker(g \otimes
\text{id}_Q) \to \Ker(f \otimes \text{id}_Q)
\to \Ker(f' \otimes \text{id}_Q) \to 0
$$
for every $R$-module $Q$.  So $f'$ is universally injective if and only if
$\Ker(f \otimes \text{id}_Q ) \subset \Ker(g \otimes
\text{id}_Q)$ for every $Q$, if and only if $g$ dominates $f$.
\end{proof}


\begin{lemma}
\label{lemma-domination-fp}
Let $f: M \to N$ and $g: M \to M'$ be maps of $R$-modules.
Then $g$ dominates $f$ if and only if for any finitely presented $R$-module
$Q$, we have $\Ker(f \otimes_R \text{id}_Q) \subset \Ker(g \otimes_R
\text{id}_Q)$.
\end{lemma}

\begin{proof}
Suppose $\Ker(f \otimes_R \text{id}_Q) \subset \Ker(g \otimes_R
\text{id}_Q)$ for all finitely presented modules $Q$.  If $Q$ is an
arbitrary module, write $Q = \colim_{i \in I} Q_i$ as a colimit of a
directed
system of finitely presented modules $Q_i$.  Then $\Ker(f \otimes_R
\text{id}_{Q_i}) \subset \Ker(g \otimes_R \text{id}_{Q_i})$ for
all $i$.  Since taking directed colimits is exact and commutes with tensor
product, it follows that $\Ker(f \otimes_R \text{id}_Q) \subset \Ker(g
\otimes_R \text{id}_Q)$.
\end{proof}


\begin{lemma}
\label{lemma-kernel-tensored-fp}
Let $M$ be an $R$-module, $P$ a finitely presented $R$-module, and $f: P
\to M$ a map.  Let $Q$ be an $R$-module and suppose $x \in \Ker(P
\otimes Q \to M \otimes Q)$.  Then there exists a finitely presented
$R$-module $P'$ and a map $f': P \to P'$ such that $f$ factors through
$f'$ and $x \in \Ker(P \otimes Q \to P' \otimes Q)$.
\end{lemma}

\begin{proof}
Write $M$ as a colimit $M = \colim_{i \in I} M_i$ of a directed system of
finitely presented modules $M_i$.  Since $P$ is finitely presented, the map $f:
P \to M$ factors through $M_j \to M$ for some $j \in I$.  Upon
tensoring by $Q$ we have a commutative diagram
$$
\xymatrix{
& M_j \otimes Q \ar[dr] & \\
P \otimes Q \ar[ur] \ar[rr] & & M \otimes Q .
}
$$
The image $y$ of $x$ in $M_j \otimes Q$ is in the kernel of $M_j \otimes Q
\to M \otimes Q$.  Since $M \otimes Q = \colim_{i \in I} (M_i
\otimes
Q)$, this means $y$ maps to $0$ in $M_{j'} \otimes Q$ for some $j' \geq j$.
Thus we may take $P' = M_{j'}$ and $f'$ to be the composite $P \to M_j
\to M_{j'}$.
\end{proof}


\begin{proposition}
\label{proposition-fp-tensor}
Let $M$ be an $R$-module.  The following are equivalent:
\begin{enumerate}
\item $M$ is finitely presented.
\item For every family $(Q_{\alpha})_{\alpha \in A}$ of $R$-modules, the
canonical map $M \otimes_R \left( \prod_{\alpha} Q_{\alpha} \right)
\to \prod_{\alpha} (M \otimes_R Q_{\alpha})$ is bijective.
\item For every $R$-module $Q$ and every set $A$, the canonical map $M
\otimes_R Q^{A} \to (M \otimes_R Q)^{A}$ is bijective.
\item For every set $A$, the canonical map $M \otimes_R R^{A} \to
M^{A}$ is bijective.
\end{enumerate}
\end{proposition}

\begin{proof}
First we prove (1) implies (2).  Choose a presentation $R^m \to R^n
\to M$ and consider the commutative diagram
$$
\xymatrix{
R^m \otimes_R (\prod_{\alpha} Q_{\alpha}) \ar[r] \ar[d]^{\cong} & R^n
\otimes_R (\prod_{\alpha} Q_{\alpha}) \ar[r] \ar[d]^{\cong} & M \otimes_R
(\prod_{\alpha} Q_{\alpha}) \ar[r] \ar[d] & 0 \\
\prod_{\alpha} (R^m \otimes_R Q_{\alpha}) \ar[r] & \prod_{\alpha} (R^n
\otimes_R Q_{\alpha}) \ar[r] & \prod_{\alpha} (M \otimes_R Q_{\alpha})
\ar[r] & 0.
}
$$
The first two vertical arrows are isomorphisms and the rows are exact.  This
implies that the map
$M \otimes_R (\prod_{\alpha} Q_{\alpha}) \to
\prod_{\alpha} ( M \otimes_R Q_{\alpha})$
is surjective and, by a diagram chase, also injective.  Hence (2) holds.

\medskip\noindent
Obviously (2) implies (3) implies (4), so it remains to prove (4) implies (1).
From Proposition \ref{proposition-fg-tensor}, if (4) holds we already know that
$M$ is finitely generated.  So we can choose a surjection $F \to M$
where $F$ is free and finite.  Let $K$ be the kernel.  We must show $K$ is
finitely generated.  For any set $A$, we have a commutative diagram
$$
\xymatrix{
& K \otimes_R R^A \ar[r] \ar[d]_{f_3} & F \otimes_R R^A \ar[r]
\ar[d]_{f_2}^{\cong} & M \otimes_R R^A \ar[r] \ar[d]_{f_1}^{\cong} & 0 \\
0 \ar[r] & K^A \ar[r] & F^A \ar[r] & M^A \ar[r] & 0 .
}
$$
The map $f_1$ is an isomorphism by assumption, the map $f_2$ is an isomorphism
since $F$ is free and finite, and the rows are exact.  A diagram chase shows
that $f_3$ is surjective, hence by Proposition \ref{proposition-fg-tensor} we
get that $K$ is finitely generated.
\end{proof}


\begin{proposition}
\label{proposition-fg-tensor}
Let $M$ be an $R$-module.  The following are equivalent:
\begin{enumerate}
\item $M$ is finitely generated.
\item For every family $(Q_{\alpha})_{\alpha \in A}$ of $R$-modules, the
canonical map $M \otimes_R \left( \prod_{\alpha} Q_{\alpha} \right)
\to \prod_{\alpha} (M \otimes_R Q_{\alpha})$ is surjective.
\item For every $R$-module $Q$ and every set $A$, the canonical map $M
\otimes_R Q^{A} \to (M \otimes_R Q)^{A}$ is surjective.
\item For every set $A$, the canonical map $M \otimes_R R^{A} \to
M^{A}$ is surjective.
\end{enumerate}
\end{proposition}

\begin{proof}
First we prove (1) implies (2).  Choose a surjection $R^n \to M$ and
consider the commutative diagram
$$
\xymatrix{
R^n \otimes_R (\prod_{\alpha} Q_{\alpha})  \ar[r]^{\cong} \ar[d] &
\prod_{\alpha} (R^n \otimes_R Q_{\alpha}) \ar[d] \\
M \otimes_R (\prod_{\alpha} Q_{\alpha})  \ar[r] & \prod_{\alpha} ( M
\otimes_R Q_{\alpha}).
}
$$
The top arrow is an isomorphism and the vertical arrows are surjections.  We
conclude that the bottom arrow is a surjection.

\medskip\noindent
Obviously (2) implies (3) implies (4), so it remains to prove (4) implies (1).
In fact for (1) to hold it suffices that the element $d = (x)_{x \in M}$ of
$M^M$ is in the image of the map $f: M \otimes_R R^{M} \to M^M$.  In
this case $d = \sum_{i = 1}^{n} f(x_i \otimes a_i)$ for some $x_i \in M$ and
$a_i \in R^M$.  If for $x \in M$ we write $p_x: M^M \to M$ for the
projection onto the $x$-th factor, then
$$
x = p_x(d) = \sum\nolimits_{i = 1}^{n} p_x(f(x_i \otimes a_i)) =
\sum\nolimits_{i=1}^{n} p_x(a_i) x_i.
$$
Thus $x_1, \ldots, x_n$ generate $M$.
\end{proof}


\begin{lemma}
\label{lemma-module-colimit-fp}
Let $R$ be a ring and let $M$ be an $R$-module.
Then $M$ is the colimit of a directed system
$(M_i, \mu_{ij})$ of $R$-modules
with all $M_i$ finitely presented $R$-modules.
\end{lemma}

\begin{proof}
Consider any finite subset $S \subset M$ and any finite
collection of relations $E$ among the elements
of $S$. So each $s \in S$ corresponds to $x_s \in M$ and
each $e \in E$ consists of a vector
of elements $f_{e, s} \in R$ such that $\sum f_{e, s} x_s = 0$.
Let $M_{S, E}$ be the cokernel of the map
$$
R^{\# E} \longrightarrow R^{\# S}, \quad
(g_e)_{e\in E} \longmapsto (\sum g_e f_{e, s})_{s\in S}.
$$
There are canonical maps $M_{S, E} \to M$.
If $S \subset S'$ and if the elements of
$E$ correspond, via this map, to relations
in $E'$, then there is an obvious map
$M_{S, E} \to M_{S', E'}$ commuting with the
maps to $M$. Let $I$ be the set of pairs
$(S, E)$ with ordering by inclusion as above.
It is clear that the colimit of this directed system is $M$.
\end{proof}


\begin{proposition}
\label{proposition-ML-tensor}
Let $M$ be an $R$-module.  The following are equivalent:
\begin{enumerate}
\item $M$ is Mittag-Leffler.
\item For every family $(Q_{\alpha})_{\alpha \in A}$ of $R$-modules, the
canonical map $M \otimes_R \left( \prod_{\alpha} Q_{\alpha} \right)
\to \prod_{\alpha} (M \otimes_R Q_{\alpha})$ is injective.
\end{enumerate}
\end{proposition}

\begin{proof}
First we prove (1) implies (2).  Suppose $M$ is Mittag-Leffler and let $x$ be
in the kernel of $M \otimes_R (\prod_{\alpha} Q_{\alpha}) \to
\prod_{\alpha} (M \otimes_R Q_{\alpha})$.  Write $M$ as a colimit $M =
\colim_{i \in I} M_i$ of a directed system of finitely presented modules
$M_i$.
 Then $M \otimes_R (\prod_{\alpha} Q_{\alpha})$ is the colimit of $M_i
\otimes_R (\prod_{\alpha} Q_{\alpha})$.  So $x$ is the image of an element
$x_i \in M_i \otimes_R (\prod_{\alpha} Q_{\alpha})$.  We must show that $x_i$
maps to $0$ in $M_j \otimes_R (\prod_{\alpha} Q_{\alpha})$ for some $j \geq
i$.  Since $M$ is Mittag-Leffler, we may choose $j \geq i$ such that $M_i
\to M_j$ and $M_i \to M$ dominate each other.  Then consider
the commutative diagram
$$
\xymatrix{
M \otimes_R (\prod_{\alpha} Q_{\alpha}) \ar[r] & \prod_{\alpha} (M
\otimes_R Q_{\alpha}) \\
M_i \otimes_R (\prod_{\alpha} Q_{\alpha}) \ar[r]^{\cong} \ar[d] \ar[u] &
\prod_{\alpha} (M_i \otimes_R Q_{\alpha}) \ar[d] \ar[u] \\
M_j \otimes_R (\prod_{\alpha} Q_{\alpha}) \ar[r]^{\cong} & \prod_{\alpha}
(M_j \otimes_R Q_{\alpha})
}
$$
whose bottom two horizontal maps are isomorphisms, according to Proposition
\ref{proposition-fp-tensor}.  Since $x_i$ maps to $0$ in $\prod_{\alpha} (M
\otimes_R Q_{\alpha})$, its image in $\prod_{\alpha} (M_i \otimes_R
Q_{\alpha})$ is in the kernel of the map $\prod_{\alpha} (M_i \otimes_R
Q_{\alpha}) \to \prod_{\alpha} (M \otimes_R Q_{\alpha})$.  But this
kernel equals the kernel of $\prod_{\alpha} (M_i \otimes_R Q_{\alpha})
\to \prod_{\alpha} (M_j \otimes_R Q_{\alpha})$ according to the
choice of $j$.  Thus $x_i$ maps to $0$ in $\prod_{\alpha} (M_j \otimes_R
Q_{\alpha})$ and hence to $0$ in $M_j \otimes_R (\prod_{\alpha} Q_{\alpha})$.

\medskip\noindent
Now suppose (2) holds. We prove $M$ satisfies formulation (1) of being
Mittag-Leffler from Proposition \ref{proposition-ML-characterization}.  Let $f:
P \to M$ be a map from a finitely presented module $P$ to $M$.  Choose
a set $B$ of representatives of the isomorphism classes of finitely presented
$R$-modules. Let $A$ be the set of pairs $(Q, x)$ where $Q \in B$ and $x \in
\Ker(P \otimes Q \to M \otimes Q)$.  For $\alpha = (Q, x) \in A$, we
write $Q_{\alpha}$ for $Q$ and $x_{\alpha}$ for $x$.  Consider the commutative
diagram
$$
\xymatrix{
M \otimes_R (\prod_{\alpha} Q_{\alpha}) \ar[r] &
\prod_{\alpha} (M \otimes_R Q_{\alpha}) \\
P \otimes_R (\prod_{\alpha} Q_{\alpha}) \ar[r]^{\cong} \ar[u] &
\prod_{\alpha} (P \otimes_R Q_{\alpha}) \ar[u] .
}
$$
The top arrow is an injection by assumption, and the bottom arrow is an
isomorphism by Proposition \ref{proposition-fp-tensor}.  Let $x \in P
\otimes_R (\prod_{\alpha} Q_{\alpha})$ be the element corresponding to
$(x_{\alpha}) \in \prod_{\alpha} (P \otimes_R Q_{\alpha})$ under this
isomorphism.  Then $x \in \Ker( P \otimes_R (\prod_{\alpha} Q_{\alpha})
\to M \otimes_R (\prod_{\alpha} Q_{\alpha}))$ since the top arrow in
the diagram is injective.  By Lemma \ref{lemma-kernel-tensored-fp}, we get a
finitely presented module $P'$ and a map $f': P \to P'$ such that $f: P
\to M$ factors through $f'$ and $x \in \Ker(P \otimes_R
(\prod_{\alpha} Q_{\alpha}) \to P' \otimes_R (\prod_{\alpha}
Q_{\alpha}))$.  We have a commutative diagram
$$
\xymatrix{
P' \otimes_R (\prod_{\alpha} Q_{\alpha}) \ar[r]^{\cong} &
\prod_{\alpha} (P' \otimes_R Q_{\alpha}) \\
P \otimes_R (\prod_{\alpha} Q_{\alpha}) \ar[r]^{\cong} \ar[u] &
\prod_{\alpha} (P \otimes_R Q_{\alpha}) \ar[u] .
}
$$
where both the top and bottom arrows are isomorphisms by Proposition
\ref{proposition-fp-tensor}.  Thus since $x$ is in the kernel of the left
vertical map, $(x_{\alpha})$ is in the kernel of the right vertical map.  This
means $x_{\alpha} \in \Ker(P \otimes_R Q_{\alpha} \to P' \otimes_R
Q_{\alpha})$ for every $\alpha \in A$.  By the definition of $A$ this means
$\Ker(P \otimes_R Q \to P' \otimes_R Q) \supset \Ker(P \otimes_R
Q \to M \otimes_R Q)$ for all finitely presented $Q$ and, since $f: P
\to M$ factors through $f': P \to P'$, actually equality holds.
 By Lemma \ref{lemma-domination-fp}, $f$ and $f'$ dominate each other.
\end{proof}
\end{document}
